Históricamente, el desarrollo de plataformas web se ha fundamentado en el modelo de \textit{Multi-Page Application} (MPA), donde cada interacción del usuario requiere una petición al servidor y la recarga completa del documento HTML, resultando ineficaces en términos de latencia y persistencia de estado.

No obstante, la evolución hacia herramientas de alta interactividad y manipulación de datos en tiempo real ha consolidado a las \textbf{Single-Page Applications} (SPA) como el estándar de la industria, adoptado por referentes como \textit{VS Code Web}, \textit{GitHub Codespaces} o las interfaces de \textit{ChatGPT}. Con esta arquitectura, el navegador descarga inicialmente un «caparazón» o \textit{shell} de la aplicación, a partir de ese momento, la comunicación con el servidor se limita exclusivamente al intercambio de datos, generalmente en formato JSON, permitiendo que la página nunca se recargue globalmente, sino que se transforme de manera dinámica. Esta arquitectura hace que la experiencia de usuario (UX) sea similar a la de las aplicaciones de escritorio, minimizando la latencia percibida y eliminando las interrupciones visuales que entorpecen el flujo de trabajo \cite{mikowski2013single}.